\section*{Chapitre \ref{ch:g1:hot-and-cold} --- Le chaud et le froid}

\begin{solution}{exo:g1:hot-and-cold:1}
Le glaçon, puis l'air de la chambre, puis l'eau du bain, puis la
soupe qui sort de la casserole.
\end{solution}

\begin{solution}{exo:g1:hot-and-cold:2}
L'eau du robinet est \emph{froide} (ou plus fraîche) à côté du
chocolat chaud, mais \emph{chaude} (ou plus chaude) à côté du jus.
\end{solution}

\begin{solution}{exo:g1:hot-and-cold:3}
Par exemple~: le Soleil, une flamme, un radiateur, ton propre corps
(ou des mains qui frottent). C'est le Soleil qui réchauffe le parc
tout entier.
\end{solution}

\begin{solution}{exo:g1:hot-and-cold:4}
Les paumes semblent chaudes contre la joue~: le frottement les a
réchauffées.
\end{solution}

\begin{solution}{exo:g1:hot-and-cold:5}
Le glaçon au soleil fond le premier~: le Soleil le réchauffe, et
c'est la chaleur qui fait fondre la glace. Celui qui est à l'ombre
reçoit moins de chaleur et dure plus longtemps.
\end{solution}

\begin{solution}{exo:g1:hot-and-cold:6}
La flaque, c'est de l'eau --- la même eau qui formait le glaçon. La
glace dure n'a pas disparu~; la chaleur de la pièce l'a fait fondre
en eau liquide.
\end{solution}

\begin{solution}{exo:g1:hot-and-cold:7}
Après toute une matinée à la même ombre, le toboggan et le banc sont
aussi froids l'un que l'autre. Le métal \emph{paraît} seulement plus
froid à la main --- le toucher n'est pas un juge équitable, comme
l'ont montré les cuillères en métal et en bois.
\end{solution}

\begin{solution}{exo:g1:hot-and-cold:8}
Toucher un bol, puis l'autre, tout de suite, avec la \emph{même}
main, et demander à un camarade de vérifier sans avoir entendu ta
réponse. La même main est nécessaire parce que le tour des trois bols
a montré que deux mains peuvent ne pas être d'accord sur une seule et
même eau.
\end{solution}

\begin{solution}{exo:g1:hot-and-cold:9}
Léa a raison. Un manteau ne fabrique pas de chaleur~: il empêche
seulement la chaleur de passer. Le bonhomme de neige n'a aucune
chaleur à garder --- en revanche, le manteau tient à l'écart
\emph{l'air chaud de la journée}, si bien que le bonhomme enveloppé
fond plus lentement et dure plus longtemps.
\end{solution}
